%@TheDoctorRAB
%
%presentation template for class slides and research presentations
%allows for citations in slide and a list at the end
%
%%%%%
%
%REFERENCES
%
%neup.bst - numbered citations in order of appearance, short author list with et al in reference section
%nsf.bst - numbered citations in order of appearance, full author list in references section
%standard.bst - citations with author last name with et al for more than 2 authors; full author list in references section
%ans.bst is for ANS only. 
%
%author = {Lastname, Firstname and Lastname, Firstname and Lastname, Firstname} for all bst formats
%bst renders the author list itself
%
%author = {{Nuclear Regulatory Commission}} if the author is an organization, institution, etc., and not people
%
%title = {{}} for all
%
%for all - use \citep{-} - [1] or (Borrelli, 2021) in the text
%standard.bst \cite{-} - Borrelli (2021) in the text
%standard.bst lists references alphabetically
%the rest list numerically
%
%
%%% citations on slides 
%
% \citep{xxxnna} where the citation should go
%
% \blfootnote{\fontsize\cite{xxxnna}\fontsize\bibentry{xxxnna}}\\ 
% before \end{frame}
% footnotes aren't working for some reason
%
% this works
% \tiny\cite{xxxnna}\tiny\hspace{0.05in}\bibentry{xxxnna}\\ 
% before \end{frame}
%%%%%


%%%%% enable tagging for accessibility
%\RequirePackage{pdfmanagement-testphase}
%\DocumentMetadata{
%    tagging=off,
%    pdfversion=2.0,
%    pdfstandard=ua-2,
%    lang=en-US,
%    debug={log=all},
%    testphase={latest,tagpdf},
%    tagging-setup={math/setup={mathml-AF,mathml-SE}}
%}
%\tagpdfsetup{math/alt/use}
%%%%%


%%%%% presentation settings
\documentclass[aspectratio=1610,pdftex,dvipsnames,compress,xcolor={dvipsnames}]{beamer}
\usetheme{Boadilla}
\usecolortheme{seahorse}
\beamertemplatenavigationsymbolsempty
\addtobeamertemplate{footnote}{\hskip -2em}{} %pushes footnote to margin
%%% font style - can add size, etc.
%https://tug.ctan.org/macros/latex/contrib/beamer/doc/beameruserguide.pdf
\setbeamerfont{title}{series=\bfseries}
\setbeamerfont{frametitle}{series=\bfseries}
\setbeamerfont{footline}{series=\bfseries}
\setbeamerfont{author}{series=\bfseries}
\setbeamerfont{institute}{series=\bfseries}
\setbeamerfont{date}{series=\bfseries}
\setbeamertemplate{page number in head/foot}[framenumber] %just gives slide number; comment out for 1/7, 2/7...
%\setbeamertemplate{caption}[numbered] %number tables and figures
%%%%%


%%%%% colors
%http://latexcolor.com/
%https://en.wikibooks.org/wiki/LaTeX/Colors#:~:text=black%2C%20blue%2C%20brown%2C%20cyan,be%20available%20on%20all%20systems.
%sam root
\definecolor{BackGround}{RGB}{255,250,240}
\definecolor{PrideGold}{RGB}{241,179,0}
\definecolor{Silver}{RGB}{165,169,180}
\definecolor{White}{RGB}{255,255,255}
\definecolor{Black}{RGB}{25,25,25}
\definecolor{Chrome}{RGB}{245,245,245}
%%% general
\definecolor{aliceblue}{rgb}{0.94, 0.97, 1.0}
\definecolor{antiquewhite}{rgb}{0.98, 0.92, 0.84}
\definecolor{lightmauve}{rgb}{0.86, 0.82, 1.0}
\definecolor{brilliantlavender}{rgb}{0.96, 0.73, 1.0}
\definecolor{brandeisblue}{rgb}{0.0, 0.44, 1.0}
\definecolor{darkmidnightblue}{rgb}{0.0, 0.2, 0.4}
\definecolor{darkchampagne}{rgb}{0.76, 0.7, 0.5}
%%% set slides
\setbeamercolor{background canvas}{bg=BackGround}
%%%
\setbeamercolor{block title}{bg=Silver,fg=Black}
\setbeamercolor{block body}{bg=Silver!20,fg=Black}
\setbeamercolor{block title alerted}{bg=Black,fg=Silver}
\setbeamercolor{block body alerted}{bg=PrideGold,fg=Black}
\setbeamercolor{alerted text}{fg=Silver}
\setbeamercolor*{block title example}{bg=PrideGold,fg=Black}
\setbeamercolor*{block body example}{bg=PrideGold!20,fg=Black}
%%%
\setbeamercolor*{palette primary}{bg=PrideGold,fg=Black}
\setbeamercolor*{palette secondary}{bg=PrideGold,fg=Black}
\setbeamercolor*{palette tertiary}{bg=PrideGold,fg=Black}
%%%
\setbeamercolor*{titlelike}{bg=PrideGold,fg=Black}
\setbeamercolor*{title}{bg=PrideGold,fg=Black}
\setbeamercolor*{item}{fg=PrideGold}
\setbeamercolor*{caption name}{fg=PrideGold}
%%%
\setbeamercolor*{sidebar}{fg=PrideGold,bg=Black}
\setbeamercolor*{title in sidebar}{fg=PrideGold}
\setbeamercolor*{author in sidebar}{fg=PrideGold}
\setbeamercolor*{section in sidebar}{fg=PrideGold}
%%%
\setbeamercolor{section in toc}{fg=Black}
\setbeamercolor{subsection in toc}{fg=Black}
%%%
\setbeamercolor{page number in head/foot}{fg=Black,bg=PrideGold}
%\setbeamercolor{footline}{bg=Black}
%%%
\setbeamercolor{bibliography entry author}{fg=Black}
\setbeamercolor{bibliography entry note}{fg=Black}
\setbeamercolor{bibliography entry title}{fg=Black}
%%%
\newcommand{\x}{\cellcolor{aliceblue}} %use to shade in table cell
\newcommand{\y}{\cellcolor{lightgray}} %use to shade in table cell
\newcommand{\z}{\cellcolor{antiquewhite}} %use to shade in table cell
\newcommand{\w}{\cellcolor{darkchampagne}} %use to shade in table cell

%%%%%


%%%%% general 
\usepackage{cmap} %math
\usepackage[pagewise]{lineno} %line numbering
\usepackage{setspace}
\usepackage{ulem} %strikethrough - do not \sout{\cite{}}
\usepackage{graphicx}
\usepackage{mypythonhighlight,verbatim}
\usepackage{filecontents}
\usepackage{footnotehyper}
\usepackage{float}
%\usepackage{subfig}
\usepackage[yyyymmdd]{datetime} %date format
\renewcommand{\dateseparator}{.}
\graphicspath{{$TEXIMG/}} %path to graphics
\setcounter{secnumdepth}{5} %set subsection to nth level
\usepackage{needspace}
\usepackage[stable,hang,flushmargin]{footmisc} %footnotes in section titles and no indent; standard.bst
\usepackage[inline]{enumitem}
\setlist[itemize]{label=\textbullet}
\usepackage{boldline}
\usepackage{makecell}
\usepackage{booktabs}
\usepackage{amssymb}
\usepackage{gensymb}
\usepackage{amsmath,nicefrac}
\usepackage{physics}
\usepackage{lscape}
\usepackage{array}
\usepackage{chngcntr}
%\usepackage{sectsty}
\usepackage{textcomp}
\usepackage{lastpage}
\usepackage{xargs} %for \newcommandx
\usepackage[colorinlistoftodos,prependcaption,textsize=tiny]{todonotes} %makes colored boxes for commenting
\usepackage{soul}
\usepackage{color}
\usepackage{marginnote}
\usepackage[figure,table]{totalcount}
\usepackage{microtype} %improves typography for pdf
\usepackage[pdftex,dvipsnames]{colortbl} %change font color
\usepackage[font=itshape]{quoting}
\usepackage[export]{adjustbox} %export allows valign in \includegraphics
%%%%%


%%%%% links and tags
\usepackage{hyperref}
\hypersetup{
    colorlinks,
    linkcolor=black,
    citecolor=black,
    urlcolor=blue,
%    pdftitle={lecture},
%    pdfauthor={R. A. Borrelli},
%    pdfkeywords={lecture},
%    pdflanguage={en-US},
%    pdfsubject={lecture},
%    pdfdisplaydoctitle=true 
} 
\usepackage[capitalise]{cleveref}
%%%%%


%%%%% tikz
\usepackage{pgf}
\usepackage{tikz} % required for drawing custom shapes
\usetikzlibrary{shapes,arrows,automata,trees}
%%%%%


%%%%% fonts
\usepackage{times}
%\renewcommand{\sfdefault}{ubuntu}
%arial - uncomment next two lines
%\usepackage{helvet}
%\renewcommand{\familydefault}{\sfdefault}
%%%%%


%%%%% references
%\usepackage[round,semicolon]{natbib} %for (Borrelli 2021; Clooney 2019) - standard.bst 
\usepackage[numbers,sort&compress]{natbib} %for [1-3] - use neup.bst; nsf.bst isn't working
\usepackage{bibentry}
%%%
\setbeamertemplate{bibliography entry title}{} %deletes linebreaks for in slide references
\setbeamertemplate{bibliography entry article}{}
\setbeamertemplate{bibliography entry title}{}
\setbeamertemplate{bibliography entry location}{}
\setbeamertemplate{bibliography entry note}{}
%%%
\setlength{\bibsep}{7pt} %sets space between references
%\renewcommand{\bibsection}{} %suppresses large 'references' heading
%\renewcommand\bibpreamble{\vspace{\baselineskip}} %sets spacing after heading if not using default references heading
%%%%%


%%%%% tables and figures
\usepackage{longtable} %need to put label at top under caption then \\ - use spacing
\usepackage{makecell}
\usepackage{tablefootnote}
\usepackage{tabularx}
\usepackage{multirow}
\usepackage{tabto} %general tabbed spacing
\usepackage{pdfpages}
\usepackage{wrapfig} %wraps figures around text
\setlength{\intextsep}{0.00mm}
\setlength{\columnsep}{1.00mm}
\usepackage[singlelinecheck=false,labelfont=bf]{caption}
\usepackage{subcaption}
\captionsetup[table]{justification=justified,skip=5pt,labelformat={default},labelsep=period,name={Table}} %sets a space after table caption
\captionsetup[figure]{justification=justified,skip=5pt,labelformat={default},labelsep=period,name={Figure}} %sets space above caption, 'figure' format
\captionsetup[wrapfigure]{justification=centering,aboveskip=0pt,belowskip=0pt,labelformat={default},labelsep=period,name={Fig.}} %sets space above caption, 'figure' format
\captionsetup[wraptable]{justification=centering,aboveskip=0pt,belowskip=0pt,labelformat={default},labelsep=period,name={Table}} %sets space above caption, 'figure' format
%%%%%


%%%%% watermark
%\usepackage[firstpage,vpos=0.63\paperheight]{draftwatermark}
%\SetWatermarkText{\shortstack{DRAFT\\do not distribute}}
%\SetWatermarkScale{0.20}
%%%%%


%%%%% cross referencing files
%\usepackage{xr} %for revisions - will cross reference from one file to here
%\externaldocument{/path/to/auxfilename} %aux file needed
%%%%%


%%%%% toc and glossaries
\usepackage[acronym,nomain,nonumberlist]{glossaries}
\makenoidxglossaries
%%%%%


%%%%% editing
\newcommand{\edit}[1]{\textcolor{blue}{#1}} %shortcut for changing font color on revised text
\newcommand{\fn}[1]{\footnote{#1}} %shortcut for footnote tag
\newcommand*\sq{\mathbin{\vcenter{\hbox{\rule{.3ex}{.3ex}}}}} %makes a small square as a separator $\sq$
%\newcommand{\sk}[1]{\sout{#1}} %shortcut for default strikethrough - do not sk through citep
\newcommand\sk{\bgroup\markoverwith{\textcolor{red}{\rule[0.5ex]{1pt}{1pt}}}\ULon} %strikethrough with red line; not in \section{}
%\st{} does strikethrough using soul package but does not like acronyms
\newcommand{\blucell}{\cellcolor{aliceblue}} %use to shade in table cell
\newcommand{\grycekk}{\cellcolor{lightgray}} %use to shade in table cell
\newcommand{\whicell}{\cellcolor{antiquewhite}} %use to shade in table cell
%%%%%


%%%%% acronyms
\newcommand{\acf}{\acrfull} %full acronym
\newcommand{\acl}{\acrlong} %long acronym
\newcommand{\acs}{\acrshort} %short acronym

\newcommand{\acfp}{\acrfullpl} %full acronym plural
\newcommand{\aclp}{\acrlongpl} %long acronym plural
\newcommand{\acsp}{\acrshortpl} %short acronym plural
%%%%%


%%%%% todonotes
\newcommandx{\cmt}[2][1=]{\todo[author=\textbf{STRUCTURE},tickmarkheight=0.15cm,linecolor=red,backgroundcolor=red!25,bordercolor=black,#1]{#2}}
\newcommandx{\con}[2][1=]{\todo[author=\textbf{CONTENT},tickmarkheight=0.15cm,linecolor=brilliantlavender,backgroundcolor=brilliantlavender,bordercolor=black,#1]{#2}}
%\newcommandx{\rab}[2][1=]{\todo[noline,author=\textbf{RAB},backgroundcolor=Plum!25,bordercolor=black,#1]{#2}}
%%%
%\newcommandx{\jon}[2][1=]{\todo[noline,author=\textbf{ATTN: Johnson},backgroundcolor=blue!25,bordercolor=black,#1]{#2}}
%\newcommandx{\han}[2][1=]{\todo[noline,author=\textbf{ATTN: Haney},backgroundcolor=OliveGreen!25,bordercolor=black,#1]{#2}}
\newcommandx{\rab}[2][1=]{\todo[author=\textbf{Food for thought},tickmarkheight=0.15cm,linecolor=black,backgroundcolor=Plum!25,bordercolor=black,#1]{#2}}
%\newcommandx{\han}[2][1=]{\todo[author=\textbf{ATTN: Haney},tickmarkheight=0.15cm,linecolor=OliveGreen,backgroundcolor=OliveGreen!25,bordercolor=OliveGreen,#1]{#2}}
%\newcommandx{\jon}[2][1=]{\todo[author=\textbf{ATTN: Johnson},tickmarkheight=0.15cm,linecolor=blue,backgroundcolor=blue!25,bordercolor=blue,#1]{#2}}
%%% highlighting 
\DeclareRobustCommand{\hlc}[1]{{\sethlcolor{LimeGreen}\hl{#1}}}
\makeatletter
    \if@todonotes@disabled
    \newcommand{\hlh}[2]{#1}
    \else
    \newcommand{\hlh}[2]{\han{#2}\hlc{#1}}
    \fi
    \makeatother

\DeclareRobustCommand{\hld}[1]{{\sethlcolor{CornflowerBlue}\hl{#1}}}
\makeatletter
    \if@todonotes@disabled
    \newcommand{\hlj}[2]{#1}
    \else
    \newcommand{\hlj}[2]{\jon{#2}\hld{#1}}
    \fi
    \makeatother

\DeclareRobustCommand{\hlf}[1]{{\sethlcolor{lightmauve}\hl{#1}}}
\makeatletter
    \if@todonotes@disabled
    \newcommand{\hlb}[2]{#1}
    \else
    \newcommand{\hlb}[2]{\rab{#2}\hlf{#1}}
    \fi
    \makeatother
%%%%%


%%%%% table alignments
\newcolumntype{L}[1]{>{\raggedright\let\newline\\\arraybackslash\hspace{0pt}}m{#1}} %uses \raggedright with m,p{} in table column
\newcolumntype{C}[1]{>{\centering\let\newline\\\arraybackslash\hspace{0pt}}m{#1}} %uses \raggedright with m,p{} in table column
\newcolumntype{R}[1]{>{\raggedleft\let\newline\\\arraybackslash\hspace{0pt}}m{#1}} %uses \raggedright with m,p{} in table column
%%%%%


%%%%% user commands
\newcommand\blfootnote[1]{%
  \begingroup
  \renewcommand\thefootnote{}\footnote{#1}%
  \addtocounter{footnote}{-1}%
  \endgroup
}


\makeatletter
\renewcommand{\@biblabel}[1]{#1.\hfill} %bibliography ordered list has numbers left flush
\makeatother


\AtBeginSection[]{
    \begin{frame}[plain]{}
         
         \vfill

         \centering
         \begin{beamercolorbox}[sep=8pt,center,shadow=true,rounded=true]{titlelike}
             \usebeamerfont{title}\insertsectionhead\par
         \end{beamercolorbox}

         \vfill

     \end{frame}
 }
%%%%%


%%%%% acronyms
\input{$ACRONYM/acronyms}
%%%%%


%%%%% linenumbering
%\linenumbers %toggle line numbers
%\pagewiselinenumbers %reset line numbers on new page
%\modulolinenumbers[1] %line numbering interval
%%%%%


%%%%% title page
\addtocounter{framenumber}{-1} %does not count the title slide in the slide count
\title[Nuclear Siting Tools]{A \acl{cgc} case study\\for\\Nuclear Siting Tools}
\author[\acs{ans} Annual Meeting -- Denver, CO]{R. A. Borrelli $\sq$ Denia Djoki{\'c}\\}
\institute[]{
        \includegraphics[width=0.25\textwidth,valign=c]{logo/university-of-idaho/nuclear-engineering/ne-logo.png}
        \quad\quad
        \includegraphics[width=0.10\textwidth,valign=c]{logo/university-of-michigan/fastest-path-zero/umich-fpz-stack.png}\\
        \vspace*{0.25in}
        \acl{ans} Annual Meeting -- Denver, Colorado\\
    }
\date{2026.06.03}
%\titlegraphic{\includegraphics[width=0.20\textwidth]{logo/i-crews/icrews-modified.png}}
%https://www.adobe.com/acrobat/online/pdf-to-ppt.html
%%%%%


%%%%%
%\logo{\includegraphics[width=0.05\textwidth]{logo/i-crews/icrews-modified.png}} %lower right on every slide
%%%%%


\begin{document}


\nobibliography* %allows \bibentry citations in the footers of frames; NOT affect putting bibliography slide at the end


%%%%% title page with no footer
{
    \setbeamertemplate{footline}{}
    \begin{frame}[plain]{}
        \titlepage
    \end{frame}
}
%%%%%


%\section{\acl{snf} across the \acs{usa}}


%\begin{frame}[plain]{}
%
%    \begin{figure}
%        \centering
%        \includegraphics[width=0.75\textwidth]{nuclear-fuel-cycle/back-end/nuclear-waste/nei-waste-map-us.png}
%%        \caption{}
%    \end{figure}
%
%    \vspace*{\fill}
%    \textbf{\tiny\cite{nei22a}\tiny\hspace{0.05in}\bibentry{nei22a}}
%\end{frame}


\section{\acl{cbs} of \acl{snf}}


\addtocounter{framenumber}{-1} 
\begin{frame}
    \frametitle{2021 \acl{caa} allocated funds for `expenses necessary for nuclear waste disposal activities \ldots'}
    \begin{enumerate}%[item-label=(\arabic*)]
        \item[]`\ldots including storage activities'
        \item[]\acs{doe} awarded \$26M to thirteen (now twelve) consortia across the nation for building capacity in \acf{cbs}
        \item[]Support more comprehensive community engagement to collaborate with stakeholders to manage \acs{snf} effectively
        \item[]Not actual siting
        \item[]Process and frameworks
    \end{enumerate}
\end{frame}


\begin{frame}
    \frametitle{Twelve consortia awarded across the country}
    \begin{figure}
        \centering
        \includegraphics[width=0.65\textwidth,alt={nationwide consortia}]{common-ground-consortium/consortia.png}
        \caption{Nationwide consortia}
    \end{figure}
    \textbf{\tiny\cite{doe23c}\tiny\hspace{0.05in}\bibentry{doe23c}}
\end{frame}


\section{The \acf{cgc}}


\addtocounter{framenumber}{-2} 
\begin{frame}
    \frametitle{\acs{cgc} spans several institutions across the intermountain west and Denia at the \acl{um}} 
    \begin{enumerate}%[item-label=(\arabic*)]
        \item[]With Nation Tribal Energy Association
        \item[]Ended in March 2026
        \item[]But we got a \acs{neu}!
        \item[]Centered on public input
        \item[]Research focus on historical lessons, types of public decision-making, and state-specific and/or cross-jurisdictional challenges in \acl{cbs} of critical infrastructure
        \item[]Providing recommendations to \acs{doe}
    \end{enumerate}

    \vspace*{\fill}
    \textbf{\tiny\cite{ara26a}\tiny\hspace{0.05in}\bibentry{ara26a}}\\
    \textbf{\tiny\cite{cfkoe25a}\tiny\hspace{0.05in}\bibentry{cfkoe25a}}\\
    \textbf{\tiny\cite{cfbor24a}\tiny\hspace{0.05in}\bibentry{cfbor24a}}
\end{frame}


\section{Motivation \& Goals}


\addtocounter{framenumber}{-1} 
\begin{frame}
    \frametitle{Lots of executive orders were issued in 2024 for nuclear energy growth}
    \begin{enumerate}%[item-label=(\arabic*)]
        \item[]400 \acs{gwe} by 2025!
        \item[]Deploying reactors at military bases
        \item[]Recycling \& reprocessing
        \item[]Hopefully pyroprocessing for Christina and me
        \item[]Increasing domestic fuel production
        \item[]July 4th milestone
    \end{enumerate}
\end{frame}


\begin{frame}
    \frametitle{All this growth really requires strong back end management policy}
    \begin{enumerate}%[item-label=(\arabic*)]
        \item[]More \acs{snf} facilities are going to be needed at the least
        \item[]Communities or localities may want help in terms of technical feasibility, local capacity to host a site
        \item[]There are some tools out there that could provide site screening information
        \item[]We critically reviewed and tested some of them within this context 
        \item[]Based on community accessibility and utility just for an \acs{snf} \acf{cis} for now
        \item[]Because we were only funded for that at the time
        \item[]Subset of a larger study for journal submission
        \item[]Just four today
    \end{enumerate}
\end{frame}


\section{\acf{sta}}


\addtocounter{framenumber}{-1} 
\begin{frame}
    \frametitle{\acs{sta} examines feasibility for advanced nuclear facility siting}
    \begin{enumerate}%[item-label=(\arabic*)]
        \item[]Population, socioeconomics, safety, regulations, etc.
        \item[]Comparative analysis for different sites
        \item[]Site Exploration option with economic, proximity, reference, safety data layers
    \end{enumerate}

    \vspace*{\fill}
    \textbf{\tiny\cite{fpz23b}\tiny\hspace{0.05in}\bibentry{fpz23b}}
\end{frame}


\begin{frame}
    \frametitle{\acs{sta} example for Bonneville County siting}

    \begin{columns}[c]

        \begin{column}{0.70\textwidth}
            \begin{figure}
                \centering
                \includegraphics[width=1\textwidth,alt={Bonneville county}]{environmental-justice-tools/bonneville-stand-exploration.png}
                \caption{User generated}
            \end{figure}
        \end{column}

        \begin{column}{0.30\textwidth}
            \begin{small}
            \begin{enumerate}%[item-label=(\arabic*)]
                \item[]Red -- protected lands
                \item[]Blue -- Electric generators
                \item[]Yellow -- Brownfield site
            \end{enumerate}
                \vspace{0.15in}
            \begin{enumerate}%[item-label=(\arabic*)]
                \item[]\acs{sta} could be useful for screening
                \item[]Could use more information on regulations 
                \item[]And more on the criteria themselves
            \end{enumerate}
            \end{small}
        \end{column}

    \end{columns}
\end{frame}


\section{\acf{tra}}


\addtocounter{framenumber}{-1} 
\begin{frame}
    \frametitle{\acs{tra} is a \acs{gis} web tool developed by \acs{onl} for transportation routing}
    \begin{enumerate}%[item-label=(\arabic*)]
        \item[]Databases for highway, rail, waterways
        \item[]Select origin and destination to generate route
        \item[]Optimize transportation under distance, preferred routes, time
        \item[]Layers are county boundaries, critical infrastructure, population, etc.
    \end{enumerate}

    \vspace*{\fill}
    \textbf{\tiny\cite{pet18b}\tiny\hspace{0.05in}\bibentry{pet18b}}
\end{frame}


\begin{frame}
    \frametitle{\acs{tra} national population density with railways}
    \begin{figure}
        \centering
        \includegraphics[width=0.85\textwidth,alt={population density with railways}]{environmental-justice-tools/tragis-rail.png}
        \caption{User generated}
    \end{figure}
\end{frame}


\begin{frame}
    \frametitle{\acs{tra} transport from Richland to Naturita}

    \begin{columns}[c]

        \begin{column}{0.70\textwidth}
            \begin{figure}
                \centering
                \includegraphics[width=1\textwidth,alt={richland to naturita transportation}]{environmental-justice-tools/tragis-naturita.png}
                \caption{User generated}
            \end{figure}
        \end{column}

        \begin{column}{0.30\textwidth}
            \begin{tiny}
            \begin{enumerate}%[item-label=(\arabic*)]
                \item[]Southern route includes Idaho
                \item[]Direct route goes through Idaho
                \item[]Node based generation 
                \item[]1300 miles \& 32 hours
            \end{enumerate}
                \vspace{0.15in}
            \begin{enumerate}%[item-label=(\arabic*)]
                \item[]Several data layers
                \item[]Multiple travel modes
                \item[]Nodes not shown beforehand; only after destination
                \item[]Communities might investigate locations of critical infrastructure, population
                \item[]Communities would not be planning transportation routes though
            \end{enumerate}
            \end{tiny}
        \end{column}

    \end{columns}
\end{frame}


\section{\acf{lit}}


\begin{frame}
    \frametitle{\acs{lit} (\acs{onl}) to evaluate areas for storage}
    \begin{enumerate}%[item-label=(\arabic*)]
        \item[]Compare up to two locations
        \item[]Three main areas of functionality --
            \begin{enumerate}%[item-label=(\arabic*)]
                \item[]Explore -- Evaluate areas of interest
                \item[]Access -- Select assessment considerations
                \item[]Compare -- Generate assessment reports
            \end{enumerate}
    \end{enumerate}

    \vspace*{\fill}
    \textbf{\tiny\cite{le23a}\tiny\hspace{0.05in}\bibentry{le23a}}
\end{frame}


\begin{frame}
    \frametitle{\acs{lit} site screening criteria for Bonneville county}

    \begin{columns}[c]

        \begin{column}{0.70\textwidth}
            \begin{figure}
                \centering
                \includegraphics[width=1\textwidth,alt={Bonneville county site screening}]{environmental-justice-tools/lit-bonneville.png}
                \caption{User generated}
            \end{figure}
        \end{column}

        \begin{column}{0.30\textwidth}
            \begin{small}
            \begin{enumerate}%[item-label=(\arabic*)]
                \item[]Layers include Critical Habitat for Endangered and Threatened Species, Coastal Barriers, Federally Designated Lands, etc.
                \item[]Energy facilities, floodplains, hazardous facilities
            \end{enumerate}
                \vspace{0.15in}
            \begin{enumerate}%[item-label=(\arabic*)]
                \item[]West is wide open
                \item[]East has been flagged under the criteria
            \end{enumerate}
            \end{small}
        \end{column}

    \end{columns}
\end{frame}


\begin{frame}
    \frametitle{\acs{lit} site screening for Bonneville county with all assessment considerations included}
    \begin{figure}
        \centering
        \includegraphics[width=0.75\textwidth,alt={Bonneville county site screening}]{environmental-justice-tools/lit-bonneville-assessment.png}
        \caption{User generated}
    \end{figure}
\end{frame}


\begin{frame}
    \frametitle{\acs{lit} comparison for Bonneville county \& Southwest Wyoming}

    \begin{columns}[c]

        \begin{column}{0.70\textwidth}
            \begin{figure}
                \centering
                \includegraphics[width=1\textwidth,alt={comparing Bonneville county and southwest Wyoming}]{environmental-justice-tools/lit-compare.png}
                \caption{User generated}
            \end{figure}
        \end{column}

        \begin{column}{0.30\textwidth}
            \begin{small}
            \begin{enumerate}%[item-label=(\arabic*)]
                \item[]Tried to select similar area
                \item[]Same assessment considerations
            \end{enumerate}
                \vspace{0.15in}
            \begin{enumerate}%[item-label=(\arabic*)]
                \item[]Wyoming offers more land due to lower population density
                \item[]\acs{lit} offers limited capabilities
                \item[]Granularity only down to counties
                \item[]Comparison feature is useful
            \end{enumerate}
            \end{small}
        \end{column}

    \end{columns}
\end{frame}


\section{\acf{str}}


\addtocounter{framenumber}{-1} 
\begin{frame}
    \frametitle{\acs{str} is a \acs{gis} decision tool for \acs{snf} transport}
    \begin{enumerate}%[item-label=(\arabic*)]
        \item[]Designed by \acs{doe}
        \item[]Generates routes by user selected options; e.g., minimizing distance, population, etc.
        \item[]Select various sites; e.g., infrastructure, police, airports, casinos, etc.
        \item[]Select origin, destination, transportation method
        \item[]Multiple transport modes
    \end{enumerate}

    \vspace*{\fill}
    \textbf{\tiny\cite{abk16a}\tiny\hspace{0.05in}\bibentry{abk16a}}
\end{frame}


\begin{frame}
    \frametitle{\acs{str} transport from Richland to Pueblo}

    \begin{columns}[c]

        \begin{column}{0.70\textwidth}
            \begin{figure}
                \centering
                \includegraphics[width=1\textwidth,alt={transport from Richland to Pueblo}]{nuclear-fuel-cycle/back-end/nuclear-waste/start-columbia-population.png}
                \caption{User generated}
            \end{figure}
        \end{column}

        \begin{column}{0.30\textwidth}
            \begin{small}
            \begin{enumerate}%[item-label=(\arabic*)]
                \item[]Heavy haul truck
                \item[]Minimize population with 800 m buffer
                \item[]3500 miles
            \end{enumerate}
                \vspace{0.15in}
            \begin{enumerate}%[item-label=(\arabic*)]
                \item[]Windy route due to buffer
                \item[]Not much detail
                \item[]Difficulty in selecting cities
                \item[]Could be used to identify critical sites on a route
            \end{enumerate}
            \end{small}
        \end{column}

    \end{columns}
\end{frame}


\section{Discussion}


\addtocounter{framenumber}{-1} 
\begin{frame}
    \frametitle{These tools have drawbacks} 
    \begin{enumerate}%[item-label=(\arabic*)]
        \item[]Some tools could not be included due to the change in funding mandate
        \item[]Could not get access to \acs{ors} -- Looked promising from literature
        \item[]A lot of redundancy across all tools studied -- Silos
        \item[]\acs{tra} \& \acs{str} are basically the same
        \item[]\acs{tra} seemed better with options and detail
        \item[]Same for \acs{sta} over \acs{lit}
        \item[]Major drawback is that the sites require registration
        \item[]Communities may be wary of privacy issues or if the tool owners may be collecting data without consent
        \item[]Cannot download the images as jpg, png
        \item[]Use of the tools in the literature was not prevalent
        \item[]More transparency on data sources should be provided
        \item[]Lack of open source accessibility
        \item[]Would be difficult for public at large use without assistance
    \end{enumerate}
\end{frame}


\begin{frame}
    \frametitle{On the other hand, how could we aid in community decisionmaking?} 
    \begin{enumerate}%[item-label=(\arabic*)]
        \item[]Site screening and transportation routes do offer use if only just to see critical sites
        \item[]Could be used for public comments
        \item[]Community could use tools to develop a go/no-protocol in a siting process
        \item[]Honestly, not that great overall
        \item[]Easy to see what fixes are needed
    \end{enumerate}
\end{frame}


\begin{frame}[allowframebreaks]
    \frametitle{Tools pros \& cons}
    \scriptsize
        \begin{longtable}{|L{0.14\linewidth}|L{0.25\linewidth}|L{0.23\linewidth}|L{0.13\linewidth}|}
            \caption{Summary of tools}
            \label{tab-tools-1}\\
            \hline
            \textbf{Tool}
            &\textbf{Use}
            &\textbf{Features}
            &\textbf{Limits}
            \\
            \hline


            \href{https://fptz.org/}{\acs{sta} site}
            & For advanced reactor siting. 
            & Site Exploration methodology could be useful for community planners for screening. Selecting criteria was easy to perform.
            & Cannot download the map. Criteria lacked sufficient explanations. Registration required.
            \\
            \hline
        

            \href{https://webtragis.ornl.gov/login}{\acs{tra} site}
            & Modeling transportation routing. Uses \acs{gis} to identify optimal means of transportation for user selected parameters.
            & Intuitive \acs{gui} allows selection of a lot of layers including population density and Tribal lands. Can combine transport routes manually. 
            & Cannot download the map as a .jpg, etc. National parks, wilderness not included in layers. Cannot just click on the map to plan route; Transportation `nodes' are not shown on the map. Registration required.
            \\
            \hline


            \href{https://lite.ornl.gov/#/}{\acs{lit} site}
            & Evaluate areas of interest for interim storage using a pre-configured assessment considerations.
            & Presents a map of screening criteria and assessment considerations down to the county level. Areas can also be compared. Descriptions of the criteria and considerations is provided.  
            & Cannot download the map as a .jpg., etc. Lacking granularity. Registration required. 
            \\
            \hline

         
            \href{https://start.energy.gov/}{\acs{str} site}
            & Geospatial decision-support tool developed for evaluating routing options and other aspects of transporting \acs{snf}.
            & Variety of data layers. Printable map options. Can select different transportation modes for one shipping route. Can select for distance, time, population buffer. Can select a wide range of infrastructure to filter.
            & Hard to select cities, etc., for destinations. Have to zoom in significantly to select a destination. Map renders slowly. 
            \\
            \hline
    \end{longtable}
\end{frame}


\section{Wrapping up}


\addtocounter{framenumber}{-1} 
\begin{frame}{What is next?} 
    \begin{enumerate}%[item-label=(\arabic*)]
        \item[]There are a few more tools to check out
        \item[]Propose a methodology for \acl{cbs} incorporating some of the tools
        \item[]For effective back-end management, community autonomy can be prioritized by providing the necessary tools for informed, independent decision-making
        \item[]Finish the journal paper
    \end{enumerate}

\end{frame}


\begin{frame}[plain]{}
    \begin{figure}
        \centering
        \includegraphics[width=0.95\textwidth]{/home/borrelli/Media/Camera/20260404_120110.jpg}
%        \caption{}
    \end{figure}
\end{frame}


\begin{frame}{References}
    \bibliographystyle{neup}
    \tiny
    \bibliography{
        araujo,
        borrelli,
        doe,
        environmental-justice-tools,
        federal,
        nuclear-disposal,
        nuclear-general,
        nuclear-storage,
        nuclear-waste,
        yucca-mountain}
\end{frame}


%%%%%%%
%\begin{frame}{}
%    \begin{columns}
%
%        \begin{column}{0.50\textwidth}
%            \begin{enumerate}[series=outerlist,topsep=0pt,itemsep=21pt,leftmargin=*,label=(\arabic*)]
%                \item[]
%                \item[]
%            \end{enumerate}
%        \end{column}
%
%        \begin{column}{0.50\textwidth}
%            \begin{enumerate}[series=outerlist,topsep=0pt,itemsep=21pt,leftmargin=*,label=(\arabic*)]
%                \item[]
%                \item[]
%            \end{enumerate}
%        \end{column}
%
%    \end{columns}
%\end{frame}

%    \begin{figure}
%        \centering
%        \includegraphics[width=0.75\textwidth]{wsc.png}
%        \caption{\acs{wsc}}
%    \end{figure}
%%%%%%%


\begin{frame}[allowframebreaks]{Acronyms}
    \printnoidxglossary
\end{frame}


\end{document}
